\chapter{Scheduling, nextTick, and Deterministic Work}

Relax has \emph{two} schedulers:

\begin{itemize}
  \item A small microtask-based job queue used by the classic VDOM layer (\texttt{src/scheduler.ts}).
  \item A more capable lane-based scheduler used by the runtime/reactive layer (\texttt{runtime/scheduler.ts}).
\end{itemize}

They serve different purposes, but they share one theme: \textbf{make UI work predictable}.

Scheduling discussions tend to get vague fast ("it’s async, it’ll happen later").
In Relax, scheduling is treated as part of the renderer’s public contract: tests depend on it, and lifecycle semantics are defined by it.

This chapter is anchored in:

\begin{itemize}
  \item \texttt{src/scheduler.ts} and \texttt{src/__tests__/scheduler.test.ts}
  \item \texttt{runtime/scheduler.ts} and \texttt{runtime/__tests__/scheduler.test.ts}
\end{itemize}

\section{The VDOM microtask job queue (\texttt{src/scheduler.ts})}

The VDOM layer uses a minimal scheduling primitive: \texttt{enqueueJob(job)}.

\subsection{Why it exists}

Several operations in Relax are intentionally deferred:

\begin{itemize}
  \item component \texttt{onMounted()} is scheduled after mount completes
  \item component \texttt{onUnmounted()} is scheduled after unmount begins
\end{itemize}

Deferring hook execution avoids re-entrancy surprises and gives the renderer a clean boundary: "render/mount now, run hooks shortly after".

\paragraph{The hidden constraint: DOM consistency and re-entrancy}

If \texttt{onMounted()} ran inside \texttt{mountDOM()}, a component could synchronously call \texttt{updateState()}, which triggers patching while the initial mount is still underway.
By scheduling hooks, Relax ensures the mount/patch/destroy primitives finish their structural work before user hook code runs.

\subsection{Data model}

\texttt{src/scheduler.ts} maintains:

\begin{itemize}
  \item \texttt{jobs}: an array of queued callbacks
  \item \texttt{isScheduled}: a flag to ensure we schedule processing only once per tick
\end{itemize}

\subsection{enqueueJob and queueMicrotask}

\texttt{enqueueJob(job)} pushes the job and calls \texttt{scheduleUpdate()}.

\texttt{scheduleUpdate()} is the key:

\begin{verbatim}
if (isScheduled) return
isScheduled = true
queueMicrotask(processJobs)
\end{verbatim}

So no matter how many jobs you enqueue, you’ll process them in a single microtask flush.

\paragraph{Batching rule}

The \texttt{isScheduled} latch is the batching mechanism:

\begin{itemize}
  \item enqueue one job \textrightarrow schedule one microtask
  \item enqueue many more jobs before the microtask runs \textrightarrow still one microtask
\end{itemize}

This makes lifecycle hooks cheap even when many components mount in the same synchronous turn.

\subsection{processJobs: ordering and async behavior}

\texttt{processJobs()} drains \texttt{jobs} in FIFO order. For each job:

\begin{itemize}
  \item it calls \texttt{job()}
  \item it wraps the result with \texttt{Promise.resolve(result).then(...)} to log async errors
\end{itemize}

Important nuance (documented in the code): \textbf{async jobs do not block the queue}.
The queue keeps running other jobs; the promise continuation runs later.

\paragraph{Why the queue doesn’t await}

The VDOM scheduler is designed for hook-style work ("start it soon, don’t stall unrelated hooks, but surface errors").
If it awaited each job, one slow async hook could hold up every other component’s hook.

This is tested in \texttt{src/__tests__/scheduler.test.ts}:

\begin{itemize}
  \item \texttt{Asynchronous jobs enqueue more jobs} expects order \texttt{[1, 3, 2]}.
\end{itemize}

That test shows:

\begin{enumerate}
  \item async job pushes \texttt{1}
  \item it awaits a microtask, so it pauses
  \item the next (sync) job runs and pushes \texttt{3}
  \item then the async job resumes and pushes \texttt{2}
\end{enumerate}

\subsection{nextTick(): what it guarantees (and what it doesn’t)}

\texttt{nextTick()} is the testing and integration helper:

\begin{verbatim}
scheduleUpdate()
return flushPromises()
\end{verbatim}

\texttt{flushPromises()} is implemented as \texttt{setTimeout(resolve)}.
That means \textbf{nextTick resolves after the microtask flush and one macrotask turn}.

The comment in \texttt{src/scheduler.ts} is explicit:

\begin{quote}
If the jobs are asynchronous, the promise will resolve before all the jobs have completed.
\end{quote}

So the contract is:

\begin{itemize}
  \item It ensures the queued jobs have been \emph{started} and synchronous effects have run.
  \item It does not await completion of async work inside a job.
\end{itemize}

\paragraph{Practical testing guidance}

	exttt{nextTick()} is a boundary for "scheduler-driven work started", not "all async completed".
If you write a hook that does work after an \texttt{await}, you may need an additional wait (or design your test so the awaited promise resolves immediately).

Tests in \texttt{src/__tests__/scheduler.test.ts} cover:

\begin{itemize}
  \item jobs run after \texttt{nextTick}
  \item jobs run in order
  \item jobs run after synchronous code
\end{itemize}

\section{The runtime lane scheduler (\texttt{runtime/scheduler.ts})}

The runtime scheduler is built for reactive workloads and larger apps.
Instead of a single FIFO queue, it manages multiple priority lanes:

\begin{itemize}
  \item \texttt{sync}
  \item \texttt{input}
  \item \texttt{default}
  \item \texttt{transition}
  \item \texttt{idle}
\end{itemize}

This scheduler is deterministic by design.
It uses stable tie-breakers, explicit budgets, and a promotion system (aging + deadlines) so low-priority work eventually runs.

This section is grounded directly in \texttt{runtime/scheduler.ts} and \texttt{runtime/__tests__/scheduler.test.ts}.

\subsection{Creating a scheduler}

The entry points:

\begin{itemize}
  \item \texttt{createScheduler(options)}
  \item \texttt{createRequestFlush(strategy)} for how flush is triggered
  \item \texttt{createBrowserScheduler(...)} convenience wrapper
\end{itemize}

\subsection{Budgets and time slicing}

A scheduler flush is run with a budget (in milliseconds). The flush loop:

\begin{itemize}
  \item picks the next task
  \item runs it
  \item checks elapsed time
  \item if budget is exceeded and work remains, schedules another flush
\end{itemize}

This behavior is tested in \texttt{runtime/__tests__/scheduler.test.ts}:

\begin{itemize}
  \item \texttt{budget enforcement: flush splits work across multiple flush calls when budget is tight}
  \item \texttt{setFrameBudget caps the budget used by scheduled flushes}
\end{itemize}

\subsection{Lane priority and deterministic ordering}

Lane priority is declared in \texttt{LANE\_PRIORITY} and selection order is:

\begin{verbatim}
sync > input > default > transition > idle
\end{verbatim}

Within a lane, tasks are sorted with deterministic tie-breakers:

\begin{enumerate}
  \item earlier deadline
  \item earlier timestamp
  \item lower task id
\end{enumerate}

Test: \texttt{orders tasks within a lane by explicit deadline (then timestamp/id)}.

In code, the deterministic ordering is implemented by sorting each lane queue using:

\begin{verbatim}
(a.deadline - b.deadline) || (a.timestamp - b.timestamp) || (a.id - b.id)
\end{verbatim}

The final tie-breaker (\texttt{id}) is the determinism anchor: it makes equal-time scheduling stable.

\subsection{Aging and promotion (starvation prevention)}

To prevent starvation, the scheduler promotes work that is overdue or old.

The implementation:

\begin{itemize}
  \item defines \texttt{AGING\_MS = 250}
  \item in \texttt{pickNextTask()}, checks queue heads in weaker lanes
  \item if a head task is overdue (deadline reached) or aged, it is promoted one lane upward
\end{itemize}

Tests:

\begin{itemize}
  \item \texttt{ages/pushes starving tasks upwards (idle -> transition)}
  \item \texttt{promotes a task when its deadline is reached (even if it is not old enough)}
  \item \texttt{starvation prevention: repeated promotions eventually let idle work run even under continuous input load}
\end{itemize}

This is the scheduler equivalent of "don’t let background work get stuck forever".

\subsection{withBudget(): integrating user work}

The scheduler also offers \texttt{withBudget(budgetMs, fn)}.
It executes user code and, if that work consumes the budget while tasks remain, schedules continuation.

Test: \texttt{withBudget schedules continuation when user code exhausts budget and tasks remain}.

\section{How these schedulers relate to the book}

In Part II and Part IV we saw that the VDOM layer uses \texttt{enqueueJob} to schedule lifecycle hooks.

In later parts of the book (signals, blocks, hydration), we’ll lean more on the runtime scheduler.
The runtime lane model is a foundation for:

\begin{itemize}
  \item fair scheduling of effects
  \item batching updates
  \item making performance tunable (frame budgets, flush strategy)
\end{itemize}

\section{Takeaways}

\begin{itemize}
  \item \texttt{src/scheduler.ts} is tiny but intentional: microtask batching + safe error logging.
  \item \texttt{nextTick()} is a boundary for tests and integration, but does not await async job completion.
  \item \texttt{runtime/scheduler.ts} adds lanes, deadlines, aging, and budgets for scalable reactive workloads.
\end{itemize}

\section{Walking the tests: what they prove}

You can treat these as the scheduler’s executable specification.

\subsection{VDOM scheduler tests (\texttt{src/__tests__/scheduler.test.ts})}

\begin{description}
  \item[Enqueued jobs run after nextTick] Proves \texttt{enqueueJob()} is deferred and does not run synchronously.
  \item[Enqueued jobs run in order] Proves FIFO execution order in \texttt{processJobs()}.
  \item[Jobs run after synchronous code] Proves microtasks run after the current synchronous call stack.
  \item[Asynchronous jobs enqueue more jobs] Proves the "async doesn’t block" rule: a job that yields allows later jobs to run, producing \texttt{[1, 3, 2]}.
\end{description}

\subsection{Runtime scheduler tests (\texttt{runtime/__tests__/scheduler.test.ts})}

The runtime scheduler suite uses a manual clock (\texttt{now()}) and a manual flush trigger (\texttt{requestFlush(cb)} pushes \texttt{cb} into an array).
That technique eliminates timing flakiness and makes time-based scheduling deterministic.

Key tests and the guarantees they enforce:

\begin{description}
  \item[runs tasks by lane priority] Proves lane selection order and that \texttt{sync} runs immediately.
  \item[budget enforcement / setFrameBudget] Proves time slicing across multiple flush calls when the budget is tight.
  \item[ages/pushes starving tasks upwards] Proves aging-based promotion (\texttt{AGING\_MS = 250}).
  \item[promotes a task when its deadline is reached] Proves deadline-based promotion even when a task isn’t old.
  \item[withBudget schedules continuation] Proves user code that consumes budget can trigger a continuation flush if tasks remain.
  \item[createRequestFlush(messageChannel/raf)] Proves flush strategy selection when platform APIs exist.
\end{description}

\section{Online resources}

\begin{itemize}
  \item MDN: \href{https://developer.mozilla.org/en-US/docs/Web/API/queueMicrotask}{\texttt{queueMicrotask()}}
  \item MDN: \href{https://developer.mozilla.org/en-US/docs/Web/JavaScript/Event_loop}{JavaScript event loop (microtasks vs. macrotasks)}
  \item WHATWG HTML: \href{https://html.spec.whatwg.org/multipage/webappapis.html#microtask-queue}{Microtask queue (spec)}
  \item React (conceptual parallel): \href{https://react.dev/learn/queueing-a-series-of-state-updates}{Queued updates and batching}
  \item MDN (conceptual parallel): \href{https://developer.mozilla.org/en-US/docs/Web/API/Background_Tasks_API}{Cooperative background work}
\end{itemize}
